\addcontentsline{toc}{subsubsection}{Solution 55. The Reversible Difference}

\begin{dlsummary}[title={\bfseries Solution 55. The Reversible Difference},breakable]

\begin{proof}
The only solutions are
\[
\boxed{
f(n)=n^2
\quad\text{and}\quad
f(n)=\frac{n(n+1)}2
\qquad(n\in\mathbb Z).
}
\]

\paragraph{1. Reverse the step to force equality.}

Replace \(m\) by \(m+3n\) and \(n\) by \(-n\). The given inequality becomes
\[
f(m)+3f(m+2n)
\ge3f(m+n)+f(m+3n),
\]
which is precisely the reverse inequality.

Therefore
\[
\boxed{
f(m+3n)-3f(m+2n)+3f(m+n)-f(m)=0.
} \tag{1}
\]

\paragraph{2. Recover a quadratic formula.}

Take \(n=1\) in (1), and define
\[
d_m=f(m+2)-2f(m+1)+f(m).
\]
Equation (1) says
\[
d_{m+1}=d_m.
\]
Thus all second differences equal a fixed integer \(k\).

Using \(f(0)=0\) and \(f(1)=1\), the resulting recurrence gives
\[
\boxed{
f(n)=n+k\frac{n(n-1)}2
\qquad(n\in\mathbb Z).
} \tag{2}
\]
For completeness, the right-hand side has the same second differences and initial values as \(f\); the recurrence determines values both forwards and backwards.

\paragraph{3. Use nonnegativity to restrict \(k\).}

At \(n=-1\), formula (2) gives
\[
f(-1)=k-1.
\]
Since \(f(-1)\ge0\),
\[
\boxed{k\ge1.}
\]

\paragraph{4. Use the composition bound to finish.}

We have
\[
f(2)=k+2,
\]
and hence
\[
f(f(2))
=k+2+\frac{k(k+1)(k+2)}2.
\]
If \(k\ge3\), then
\[
f(f(2))\ge5+\frac{3\cdot4\cdot5}{2}=35>16,
\]
a contradiction. Consequently,
\[
k\in\{1,2\}.
\]

Substituting these values into (2) gives
\[
k=1:\quad f(n)=\frac{n(n+1)}2,
\qquad
k=2:\quad f(n)=n^2.
\]

\paragraph{5. Verify both solutions.}

Both formulas give nonnegative integers for every integer \(n\), with the required values at \(0\) and \(1\). Every quadratic satisfies (1), so the original inequality holds with equality.

Finally,
\[
f(f(2))=
\begin{cases}
6,&f(n)=n(n+1)/2,\\
16,&f(n)=n^2.
\end{cases}
\]
Both satisfy the composition bound.
\end{proof}

\paragraph{Takeaways.}
\emph{No takeaways were provided in the source draft.}

\end{dlsummary}
